\documentclass{article}
\usepackage{amsmath}

\begin{document}

\textbf{Proposition:} If $A$ and $B$ are events such that $A \subseteq B$, then $P(A) \leq P(B)$.

\textbf{Proof:}

Let $A$, $B$ be events with $A \subseteq B$. Then,

\begin{equation*}
B=A \cup (B \cap A^c)
\end{equation*}

Since $A$ and $B \cap A^c$ are mutually exclusive, it follows that

\begin{align*}
P(B) &= P(A \cup (B \cap A^c)) \\
&= P(A) + P(B \cap A^c) \\
\end{align*}

We know that $P(B \cap A^c) \geq 0$. Hence,
by adding $P(A)$ on both sides, we obtain $P(B) \geq P(A)$.

Therefore, if $A \subseteq B$, then $P(A) \leq P(B)$.

\end{document}
